\documentclass[book.tex]{subfiles}

\begin{document}

\chapter{Semantics for Differentiable Programs}

\label{chap:grammar}
\noindent\rule{11cm}{0.4pt} \\
\textbf{Prerequisites:} \autoref{chap:autodiff}, \autoref{chap:hindley_milner}  \\
\textbf{Difficulty Level:} *** \\
\noindent\rule{11cm}{0.4pt}
\vspace{5mm}

While we have studied semantics for simply typed and polymoriphically typed languages in the last few chapters, we have not yet seen how to model the semantics of a differentiable programming language. In this chapter, we present $\lambda_S$, a system of semantics for a differentiable programming language. 

\section{Physika is a Functional Language}

Throughout this book, we have given Physika code in what looks like imperative style. Most numerical programming and deep learning texts provide examples in imperative code, so allowing Physika to follow the same style simplifed exposition considerably. However, implementing automatic differentiation in an imperative language causes considerable additional difficulty since the values of variables can change during execution.

For this reason, underneath the hood, Physika is actually a fully functional programming language. Let's work through a code example. Consider the imperative code

\begin{lstlisting}
sum = 0
for i in 100:
  sum += i
\end{lstlisting}

Physika converts this loop into code that looks roughly like
\begin{lstlisting}[language=python]
sum = 0
sum0 = sum + 0
sum1 = sum0 + 1
sum2 = sum1 + 2
...
\end{lstlisting}
That is, the loop is fully unrolled and each update to a variable actually creates an entirely new variable.

As another example, consider
\begin{lstlisting}[language=python]
setting = 3
def g():
  setting += 1
  setting
print(g())
\end{lstlisting}
This code is unrolled into
\begin{lstlisting}[language=python]
setting = 3
setting = setting + 1
print(setting)
\end{lstlisting}
This unrolling operation allows for every Physika control structure to be unrolled into a simple trace that can be automatically differentiated.

\section{Converting to Traces}

The core idea of how we assign semantics to Physika programs is that we simplify Physika programs into linear traces. The idea is that a complex program which may involve control structures is transformed into a simple list of instructions. This simpler architecture can then be directly backpropagated through in order to generate the needed gradients. Note that such a simplification tranform is always possible since any program must at the end of the transform down to simple instructions that can be run on assembly.

We define a sub-language, the Physika trace language that is a restricted subset of Physika that has no conditional statements, no function definitions, no function calls, and no reverse mode differentiation. We define a series of transformations that transform an arbitrary Physika program into a trace language call. After the simplification transformations, Physika will be reduced to a core language that satisfies the following grammar.

\begin{align}
    C ::= & x  \\
    &| r \in \mathbb{R} \\
    &| C + D \\
    &| (C, D) \\
    &| C[i]   \\
    &| C(D)
\end{align}

\section{Denotational Semantics}

Denotational semantics provide a mechanism to turn a Physika program into a function on a suitable mathematical space.

\begin{align}
    \llbracket R \rrbracket &= \mathbb{R} \\
    \llbracket T \times U \rrbracket &= \llbracket T \rrbracket \times \llbracket U \rrbracket \\
    \llbracket M + N \rrbracket &= \llbracket M \rrbracket + \llbracket N \rrbracket \\
    \llbracket (M, N) \rrbracket &= (\llbracket M \rrbracket, \llbracket N \rrbracket)
\end{align}

The above fragment provides a partial listing of the denotational semantics for Physika programs that demonstrates how arithmetic operations can be mapped to corresponding operations on the real numbers.

%\textcolor{red}{TODO: the $\lambda_S$ semantics use a cartesian monoidal category. Should we go ahead and add a section on category theory to part 1a? Some of the more modern approaches to the path integral use category theory extensively so this background material may prove useful in part 5 as well.}

%\textcolor{red}{TODO: Clarke derivatives are crucial for this construction. These are a fairly exotic concept that I haven't seen elsewhere in the literature}

%\textcolor{red}{TODO: Use simple semantics from \url{https://arxiv.org/abs/1911.04523}}

%\section{Syntax}

%\textcolor{red}{TODO: Fill out syntax section based on $\lambda_S$ syntax}

%\section{Semantics}
%
%\textcolor{red}{TODO: Fill out semantics section based on $\lambda_S$ semantics}

\end{document}